\subsection{Fazit}
\label{subsec:fazit}

Über alle vier Pipeline-Stufen zeigt sich ein klares Muster: Shared Memory
Tiling lohnt sich nur dort, wo redundante Global-Memory-Zugriffe stark ins
Gewicht fallen, etwa beim Gaussian Blur mit großem Radius
(Abschnitt~\ref{sec:results_gaussian}). Bei den kleinen Nachbarschaften von
Sobel, NMS und Hysteresis überwiegt dagegen der Synchronisations-Overhead,
sodass die naive bzw. unrolled Variante schneller bleibt. Pinned Memory
verbessert durchgängig den D$\to$H-Transfer, ohne den Kernel zu verändern.
Den größten Effekt erzielt die Fused Pipeline
(Abschnitt~\ref{sec:results_fused}), die durch den Wegfall der
Zwischentransfers die Gesamtlaufzeit um über 80\,\% reduziert.

Aus der Sicht von Matthias Stefan zeigte dieses Praktikum insbesondere die
Herausforderungen sauberen Performance-Messens in der Praxis auf: Ohne
kontrollierte Laborbedingungen, etwa hinsichtlich konstanter GPU-Temperatur
oder isolierter Hintergrundprozesse (vgl.
Abschnitt~\ref{subsec:messmethodologie}), unterliegen Messwerte einer gewissen
Schwankungsbreite. Daraus ergibt sich eine grundsätzlichere methodische
Frage: Inwieweit sind unter Laborbedingungen gewonnene Messwerte
aussagekräftig, wenn die Zielumgebung, insbesondere im Consumer-Bereich,
diese Bedingungen ohnehin nie erfüllt? Realitätsnahe, durch
Hintergrundlast beeinflusste Messbedingungen könnten den späteren
Einsatzkontext unter Umständen treffender abbilden als ein steriles
Setup. Die im Praktikum gewonnenen Erkenntnisse zu Speicherzugriffsmustern
und Transfer-Overhead sind zudem auf angrenzende Bereiche wie
3D-Rendering und Vulkan übertragbar.

Aus der Sicht von Thomas Jurczyk <<<TODO>>>

Als Ausblick bleibt festzuhalten, dass sich alle Messungen auf ein einzelnes,
wenn auch großes Bild (5120$\times$2880) beschränken. Effekte wie
Speicherbandbreitenlimitierung oder Occupancy-Grenzen dürften bei deutlich
größeren oder kleineren Bildgrößen anders bzw. stärker ausfallen, sodass eine
systematische Analyse über mehrere Bildgrößen hinweg ein sinnvoller nächster
Schritt wäre.